\documentclass[../main.tex]{subfiles}
\begin{document}

Phụ lục này hướng dẫn cài đặt, cấu hình và triển khai hệ thống, giúp người đọc có thể tái
dựng lại môi trường chạy của đồ án.

\section{Yêu cầu môi trường}

Để cài đặt và chạy hệ thống, cần chuẩn bị các thành phần trong Bảng~\ref{tab:moi-truong}.

\begin{table}[H]
\centering
\small
\caption{Yêu cầu môi trường cài đặt}
\label{tab:moi-truong}
\begin{tabular}{|p{4.6cm}|p{9cm}|}
\hline
\textbf{Thành phần} & \textbf{Vai trò} \\ \hline
Node.js (bản LTS) và npm & Môi trường chạy và quản lý gói cho Mini App, trang quản trị và Cloud Functions \\ \hline
Zalo Mini App CLI (\texttt{zmp-cli}) & Chạy thử và triển khai Mini App lên nền tảng Zalo \\ \hline
Firebase CLI & Triển khai Cloud Functions, quy tắc Firestore và trang quản trị \\ \hline
Tài khoản Zalo Developer & Đăng ký và quản lý Mini App \\ \hline
Dự án Firebase & Cung cấp Firestore, Storage và Cloud Functions \\ \hline
\end{tabular}
\end{table}

\section{Cấu trúc mã nguồn}

Mã nguồn được tổ chức thành ba phần chính cùng các tệp cấu hình, như Danh sách mã nguồn~\ref{lst:cautruc}.

\begin{lstlisting}[caption={Cấu trúc thư mục mã nguồn},captionpos=b,label=lst:cautruc]
.
|- src/             # Zalo Mini App (React + TypeScript + Jotai)
|- admin/           # Trang quan tri (React + Ant Design)
|- functions/       # Cloud Functions (Node.js)
|- app-config.json  # Cau hinh Mini App
|- firestore.rules  # Quy tac bao mat Firestore
|- firebase.json    # Cau hinh Firebase CLI
\end{lstlisting}

\section{Cấu hình biến môi trường}

Các khoá bí mật và cấu hình được đặt trong các tệp môi trường và \textbf{không} đưa vào mã
nguồn công khai. Phần chạy phía trình duyệt sử dụng tiền tố \texttt{VITE\_}, còn các khoá nhạy
cảm được đặt trong \texttt{functions/.env} ở phía máy chủ. Danh sách mã nguồn~\ref{lst:env} minh hoạ cấu
trúc tệp môi trường (các giá trị chỉ là ví dụ mẫu).

\begin{lstlisting}[caption={Cấu trúc tệp môi trường (giá trị mẫu)},captionpos=b,label=lst:env]
# .env  (Mini App - phia trinh duyet)
VITE_FIREBASE_API_KEY=your_api_key
VITE_FIREBASE_PROJECT_ID=your_project_id
VITE_BYPASS_MSSV=20225413

# functions/.env  (phia may chu)
GROQ_API_KEY=your_groq_key
EMAILJS_SERVICE_ID=your_service_id
\end{lstlisting}

\section{Chạy ở môi trường phát triển}

Danh sách mã nguồn~\ref{lst:dev} trình bày các lệnh chạy thử trên máy cá nhân. Lưu ý Mini App được chạy
bằng \texttt{zmp start} (không dùng \texttt{vite} trực tiếp do cơ chế đóng gói riêng của
nền tảng).

\begin{lstlisting}[caption={Lệnh chạy môi trường phát triển},captionpos=b,label=lst:dev]
# Cai dat phu thuoc
npm install

# Chay Mini App
zmp start

# Chay trang quan tri
cd admin && npm install && npm run dev

# Kiem tra kieu TypeScript (khong xuat ban dich)
npx tsc --noEmit
\end{lstlisting}

\section{Triển khai}

Sau khi kiểm thử, hệ thống được triển khai theo Danh sách mã nguồn~\ref{lst:deploy}. Mini App sau khi
\texttt{zmp deploy} sẽ trải qua quy trình kiểm duyệt của Zalo trước khi đến tay người dùng;
trang quản trị được phục vụ trên Firebase Hosting; còn Cloud Functions chạy tại khu vực
\texttt{asia-southeast1}.

\begin{lstlisting}[caption={Lệnh triển khai},captionpos=b,label=lst:deploy]
# Trien khai Mini App len nen tang Zalo
zmp deploy

# Trien khai Cloud Functions, rules, indexes va trang quan tri
firebase deploy --only functions,firestore:rules,firestore:indexes,hosting
\end{lstlisting}

\section{Khởi tạo và phân quyền}

Tài khoản đăng nhập lần đầu mặc định mang vai trò sinh viên. Việc nâng quyền lên giảng viên
hoặc quản trị viên chỉ được thực hiện ở phía máy chủ (thông qua hàm Cloud Function tương ứng
hoặc mã mời giảng viên), nhằm ngăn việc tự nâng quyền từ phía thiết bị. Sau khi được nâng
quyền, người dùng có thể truy cập các chức năng tương ứng với vai trò như đã đặc tả ở
mục~\ref{sec:dac-ta}.

\end{document}
